\chapter{Conclusions}
\label{ch:conclusions}

This thesis presented a FRI-based polynomial commitment scheme over the
BabyBear field, developed from its mathematical foundations to a
complete, working implementation. Starting from finite fields and
Reed-Solomon codes, we described the interactive oracle proof model and
its compilation through the BCS transformation, the FRI protocol and its
soundness, and the quotient technique that turns FRI into a full
evaluation proof. Every component was implemented from scratch in
Python, and the scheme was validated by a test suite covering both the
low-degree test and the evaluation proof.

The result is a compact and readable realisation of a transparent,
hash-based commitment scheme, in which each part of the construction can
be followed directly in the code.

\section{Future Work}
\label{sec:future_work}

Several natural extensions remain. The proof-of-work grinding step,
prepared for in the parameters but left as a stub, would raise the
security level at negligible verifier cost. Beyond this, the polynomial
commitment scheme is the key ingredient of a full STARK: adding an
algebraic intermediate representation and the constraint machinery on
top of it would turn the present system into a complete proof system for
general computations. Finally, porting the implementation to a compiled
language and adding the standard batching and low-level optimisations
would bring its performance closer to that of production systems.